\section{Backward demands replace induction on completed objects}
\label{sec:backward-demands}

A recurring obstruction in an arithmetic treatment of analysis is an
induction hypothesis of the form ``the object at stage $n$ is valid at every
precision.''  For example, validity of a real name is an unbounded universal
statement.  The same issue occurs when one proves all field laws at every
stage of a tower of algebraic extensions.  Quantifier-free induction does
not permit these statements as induction formulas.

The existing packets already provide bounded Booleanization, finite loops,
and composition of observation contracts
\cite{repo-contracts,repo-embedding}.  Moreover, the Platonic-solids packet
uses a backward sequence of coefficient budgets for its field-tower
argument \cite{repo-platonic}.  The contribution here is a single explicit
rule covering nonstationary and nonmonotone demands, a construction of
parameter-changing recursions, and a checked Lean realization of their
finite scheduling arguments.  This is a generalization and reusable
formalization of that method, not a claim that backward budgeting itself
is new.

\subsection{The exact finite transfer theorem}

All functions in this subsection are fixed primitive recursive functions;
the notation $\vec a$ denotes their ordinary numerical parameters.  Let
$C(n,k,\vec a)$ be a bounded formula: it describes a finite observation of
the object at stage $n$, with demand $k$.  It can, for example, assert all
the required inequalities with indices at most $k$.  It is not required
to be monotone in $k$.  Let $D(n,k,\vec a)$ be a demand transformer.  A local
step contract has the shape
\begin{equation}\label{eq:backward-local-step}
 C(n,D(n,k,\vec a),\vec a)
       \ \longrightarrow\ C(n+1,k,\vec a).
\end{equation}
The point of this contract is that its premise asks for one specified
finite observation, instead of asking that the predecessor object satisfy
its specification at every precision.

Fix a terminal stage $N$ and a terminal demand $K$.  Define the backward
array by ordinary primitive recursion:
\begin{align}
 b(0;N,K,\vec a)&=K,\nonumber\\
 b(j+1;N,K,\vec a)
   &=D\bigl(N\mathbin{\dotminus}(j+1),b(j;N,K,\vec a),\vec a\bigr),
   \label{eq:backward-array}\\
 h(t;N,K,\vec a)&=b(N\mathbin{\dotminus}t;N,K,\vec a).
   \label{eq:backward-horizon}
\end{align}
Here $\dotminus$ is truncated subtraction.  The values with $j>N$ or
$t>N$ merely make the displayed functions total.  The proof uses only
the intended range.  Direct substitution gives
\begin{equation}\label{eq:backward-schedule-laws}
 h(N)=K,\qquad h(t)=D(t,h(t+1),\vec a)\quad(t<N).
\end{equation}
For the second equation use
$N\dotminus t=(N\dotminus(t+1))+1$ and
$N\dotminus((N\dotminus(t+1))+1)=t$, valid when $t<N$.
These are elementary numerical equations proved by quantifier-free
induction; no backwards induction principle is being added.

\begin{theorem}[Backward-observation transfer]
\label{thm:backward-transfer}
Let $\Gamma(\vec a,N,K)$ be a finite context in which the stage variable
does not occur freely.  Suppose that \ipra\ derives
\[
 \Gamma\vdash C(0,h(0;N,K,\vec a),\vec a)
\]
and, with $n,r$ free,
\[
 \Gamma,\ n<N,\ C(n,D(n,r,\vec a),\vec a)
       \vdash C(n+1,r,\vec a).
\]
Then \ipra\ derives
\[
                 \Gamma\vdash C(N,K,\vec a).
\]
Only the single initial observation at demand $h(0)$ is required.  Neither
$D$ nor $C$ needs a monotonicity hypothesis.
\end{theorem}

\begin{proof}
Fix $\vec a,N,K$ throughout the argument.  Booleanize the fixed bounded
formula $C$: let $c(n,k,\vec a)$ be its primitive recursive characteristic
function, with an \ipra\ derivation of
$C(n,k,\vec a)\leftrightarrow c(n,k,\vec a)=1$.
Use the induction formula
\begin{equation}\label{eq:backward-invariant}
 J(t):\quad
 t\le N\ \longrightarrow\ c(t,h(t;N,K,\vec a),\vec a)=1.
\end{equation}
The numerical comparison $t\le N$ is also a primitive recursive test, so
$J$ is quantifier-free in the language of \ipra.

The base follows from the supplied initial derivation and $0\le N$.
For the successor step assume $J(t)$ and $t+1\le N$.  Then $t<N$ and
$t\le N$, hence the induction hypothesis gives $C(t,h(t),\vec a)$.
By \eqref{eq:backward-schedule-laws}, this is exactly
$C(t,D(t,h(t+1),\vec a),\vec a)$.  Instantiate the supplied local step at
$n=t$ and $r=h(t+1)$.  It gives $C(t+1,h(t+1),\vec a)$, and
Booleanization yields the consequent of $J(t+1)$.  This argument proves
the implication defining $J(t+1)$; it asserts nothing about out-of-range
stages.

Quantifier-free induction proves $\forall t\,J(t)$.  Instantiate at $N$,
use $h(N)=K$, and convert the Boolean equation back to $C(N,K,\vec a)$.
The context $\Gamma$ remains fixed during induction.  Its formulas need
not themselves be decidable or bounded.
\end{proof}

\begin{corollary}[Observational induction]
\label{cor:observational-induction}
Suppose the context $\Gamma(\vec a)$ is independent of $N,K$, \ipra\
derives $\Gamma\vdash C(0,k,\vec a)$ for arbitrary $k$, and it derives
\eqref{eq:backward-local-step} under $\Gamma$.  Then it derives
\[
       \Gamma\vdash\forall n\,\forall k\ C(n,k,\vec a).
\]
\end{corollary}

\begin{proof}
Fix $N,K$, substitute $h(0;N,K,\vec a)$ for $k$ in the base derivation,
and apply Theorem~\ref{thm:backward-transfer}.  Universal introduction over
$N,K$ gives the result.  The induction formula was
\eqref{eq:backward-invariant}; it was not $\forall k\,C(t,k,\vec a)$.
\end{proof}

For a constant precision loss $c$, namely $D(n,k)=k+c$, the schedule is
especially simple:
\begin{equation}\label{eq:constant-loss-budget}
                  h(t)=K+c(N\dotminus t).
\end{equation}
If $D(n,k)=k+d(n)$, then $h(0)=K+\sum_{n<N}d(n)$.  A nonlinear $D$ is
handled by the same primitive recursion \eqref{eq:backward-array}; no
closed formula or elementary growth bound is needed.  This remains a
fixed primitive recursive definition uniformly in the numerical inputs
$N,K,\vec a$.  It is not a universal evaluator for variable program syntax.

\subsection{Parameter-changing recursion without hidden strength}

Even defining an iterated represented object can conceal a difficulty.
The equations
\begin{equation}\label{eq:parameter-changing-recursion}
 x_0(k)=b_0(k),\qquad
 x_{n+1}(k)=F(n,k,x_n(D(n,k)))
\end{equation}
change the parameter $k$ in the recursive call.  Therefore one should not
simply declare that they are an instance of the primitive-recursion
constructor with $k$ held fixed.

\begin{theorem}[Compilation of a predetermined demand recurrence]
\label{thm:parameter-changing-recursion}
For fixed primitive recursive $b_0,D,F$, there is a fixed primitive
recursive function $X(N,K)$ satisfying
\eqref{eq:parameter-changing-recursion}.  Its equations have \ipra\ proofs.
\end{theorem}

\begin{proof}
Compute the schedule $h(t;N,K)$ by
\eqref{eq:backward-array}--\eqref{eq:backward-horizon}.  Then compute
\[
 v_0=b_0(h(0)),\qquad
 v_{t+1}=F(t,h(t+1),v_t),\qquad X(N,K)=v_N.
\]
This is a composition of two ordinary primitive recursions, with $N,K$
as fixed parameters in the second one.  All data are natural-number codes.

For $N=0$, $h(0;0,K)=K$, so $X(0,K)=b_0(K)$.  For the successor
equation compare the schedule for terminal data $(N+1,K)$ with that for
$(N,D(N,K))$.  Their demands at time $N$ are both $D(N,K)$.  The same
backward recurrence therefore gives identical demands at every time
$t\le N$.  More explicitly, quantifier-free induction on $j$ proves
\[
 b(j+1;N+1,K)=b(j;N,D(N,K));
\]
the equality of the relevant truncated subtractions supplies its step.
Substitution gives equality of the two schedules at $t\le N$.
Their initial forward values agree, and quantifier-free induction on
$t\le N$ gives equality of their forward values.  The last step of the
larger computation is consequently
\[
       X(N+1,K)=F(N,K,X(N,D(N,K))).
\]
Each induction here compares numerical terms; it does not quantify over
functions or assume an unbounded family of already constructed objects.
\end{proof}

The predetermined nature of $D$ matters.  The theorem does not license
arbitrary nested recursion in which a future demand is determined by a
not-yet-computed recursive output.  Such a definition needs a separate
primitive recursive construction or bound.  A bounded trace can support
adaptive queries, as in the earlier finite-contract frameworks, but
termination of an unbounded query process is not supplied by this theorem.

\subsection{Variable finite sets of demands}

A local mathematical argument often needs several predecessor facts.
Let $A(n,j,\vec a)$ be bounded, let $m(n,j,\vec a)$ be the number of
queries, and let $q(n,j,r,\vec a)$ be the $r$th query.  Suppose the local
derivation is
\begin{equation}\label{eq:multiple-demand-step}
 \left(\forall r<m(n,j,\vec a)\;
 A(n,q(n,j,r,\vec a),\vec a)\right)
       \longrightarrow A(n+1,j,\vec a).
\end{equation}
Define
\begin{align}
 C(n,H,\vec a)&:\Longleftrightarrow
       \forall j\le H\ A(n,j,\vec a),\label{eq:prefix-observations}\\
 D(n,H,\vec a)&=
       \max_{j\le H}\ \max_{r<m(n,j,\vec a)}q(n,j,r,\vec a).
       \label{eq:multiple-demand-envelope}
\end{align}
An empty maximum is zero.  Finite maxima are primitive recursive, even
when the number of queries varies with the input and none of their indices
is monotone.

\begin{corollary}[Finite-query observational induction]
\label{cor:finite-query-induction}
Given \ipra\ derivations of $A(0,j,\vec a)$ and
\eqref{eq:multiple-demand-step} under a fixed context, there is an \ipra\
derivation of $\forall n\,\forall j\,A(n,j,\vec a)$ under that context.
\end{corollary}

\begin{proof}
For $j\le H$ and $r<m(n,j,\vec a)$, the defining maximum gives
$q(n,j,r,\vec a)\le D(n,H,\vec a)$.  Thus $C(n,D(n,H,\vec a),\vec a)$
supplies every premise in \eqref{eq:multiple-demand-step}.  It yields
$A(n+1,j,\vec a)$ for each $j\le H$, hence $C(n+1,H,\vec a)$.
Apply Corollary~\ref{cor:observational-induction} to the prefix predicate.

There are two bounded conjunctions here, not a hidden universal invariant.
Concretely, if $a(n,j,\vec a)$ Booleanizes $A$, then $C(n,H,\vec a)$ is
Booleanized by the product of the $H+1$ Boolean values $a(n,j,\vec a)$.
The query premise is the analogous product over $r<m(n,j,\vec a)$.
Their prefix-fold equations and the finite-maximum inequalities have
quantifier-free induction proofs.  The proof remains within \ipra\ when
$H$ and the query counts are input-dependent.
\end{proof}

Multidimensional observations require no new logical rule.  Code a tuple
of indices by a fixed primitive recursive pairing function, and put
domain and range checks into its bounded predicate.  A query can carry a
tag naming which law, coordinate, or estimate is requested.  If a local
step uses several earlier stages, replace \eqref{eq:prefix-observations}
by the bounded history predicate
\[
             C(n,H):\Longleftrightarrow
                  \forall i\le n\;\forall j\le H\ A(i,j).
\]
Use a demand envelope at least $H$ and at least every requested query
index, and require each requested stage to be at most $n$.  Old history
entries persist because the envelope is at least $H$; the local step
establishes the new entries.  This proves the same transfer rule for a
finite directed dependency graph ordered by its stage index.  Cyclic
dependencies are not justified by this argument.

\subsection{Worked application: nonlinear iteration of a real name}
\label{subsec:nonlinear-name-iteration}

This example exhibits the computational content of the rule.  It avoids
both an induction on equality of real numbers and the assumption that
an arbitrary nested approximation program is primitive recursive.

Fix a rational $a$ with $0\le a\le4$ and a fixed primitive recursive
rational family $u(k)$.  Suppose its regularity observations are
\[
 \operatorname{Obs}_u(H):\Longleftrightarrow
  \forall i,j\le H\quad
       |u(i)-u(j)|\le2^{-i}+2^{-j},
\]
including positivity checks for all denominators used.  Put
\[
 p(r)=\min(1,\max(0,r)),\qquad T_a(r)=ar(1-r),
\]
and define the output family directly by the finite loop
\begin{equation}\label{eq:nonlinear-name-definition}
                v(n,k)=T_a^{\,n}\bigl(p(u(k+2n))\bigr).
\end{equation}
For each input $n,k$, this evaluates $u$ once and performs $n$ exact
rational updates.  Therefore it is a fixed primitive recursive family,
uniformly in $n,k,a$ and the numerical parameters of the fixed evaluator
$u$.  Rational arithmetic may be implemented by signed numerators and
positive denominators; reduction to lowest terms is unnecessary.

Three rational inequalities provide all local mathematics.  First,
$p$ takes values in $[0,1]$ and is nonexpansive:
$|p(r)-p(s)|\le|r-s|$.  To see the latter, assume $r\le s$ and split at
$0,1$.  On the middle interval $p$ is the identity; on either outer
interval it is constant.  Crossing an endpoint removes a nonnegative
part of the distance, so in every case
$0\le p(s)-p(r)\le s-r$.
Second, $T_a$ maps $[0,1]$ into itself, because
\[
       0\le r(1-r)\le\frac14
       \quad\text{follows from}\quad (2r-1)^2\ge0.
\]
Third, for $r,s\in[0,1]$,
\begin{equation}\label{eq:logistic-lipschitz}
 |T_a(r)-T_a(s)|
    =a|r-s|\,|1-r-s|\le4|r-s|.
\end{equation}
All cases and inequalities concern rational numbers and can be checked
after clearing positive denominators.

The finite iterate definition also gives the exact identity
\begin{equation}\label{eq:logistic-parameter-shift}
                     v(n+1,k)=T_a(v(n,k+2)).
\end{equation}
Indeed, both sides start from $p(u(k+2n+2))$ and apply the same map
$n+1$ times.  This identity can equivalently be obtained by
Theorem~\ref{thm:parameter-changing-recursion} with demand $D(n,k)=k+2$.

\begin{theorem}[Finite validity of the nonlinear iterates]
\label{thm:nonlinear-name-validity}
From $0\le a\le4$ and $\operatorname{Obs}_u(K+2N)$, \ipra\ proves
\[
 \begin{gathered}
      \forall k\le K\quad 0\le v(N,k)\le1,\\
      \forall k,\ell\le K\quad
       |v(N,k)-v(N,\ell)|\le2^{-k}+2^{-\ell}.
 \end{gathered}
\]
In particular, a globally regular input family gives a globally regular
output family at every numerical stage $N$.
\end{theorem}

\begin{proof}
Let $C(n,H)$ be the conjunction of the displayed range and regularity
assertions with $n,H$ in place of $N,K$.  It is bounded.  At $n=0$,
$v(0,k)=p(u(k))$; the range assertion and nonexpansiveness of $p$ give
$C(0,K+2N)$ from the input observation.

Assume $C(n,H+2)$.  For $k\le H$, the input $v(n,k+2)$ lies in $[0,1]$,
so \eqref{eq:logistic-parameter-shift} gives the next range assertion.
For $k,\ell\le H$, the same identity and
\eqref{eq:logistic-lipschitz} give
\[
 \begin{split}
 |v(n+1,k)-v(n+1,\ell)|
 &\le4|v(n,k+2)-v(n,\ell+2)|\\
 &\le4\bigl(2^{-(k+2)}+2^{-(\ell+2)}\bigr)\\
 &=2^{-k}+2^{-\ell}.
 \end{split}
\]
Thus $C(n,H+2)\to C(n+1,H)$ is the local finite contract.
The backward schedule is $h(t)=K+2(N-t)$ for $t\le N$.
Apply Theorem~\ref{thm:backward-transfer}.  Its induction formula is the
Boolean equation for $C(t,K+2(N-t))$, guarded by $t\le N$.
Only after this finite argument is complete do we universally quantify
over $N,K$ when the input is globally regular.
\end{proof}

In the ordinary external interpretation of regular names, if $u$ names
$x$, then $v(N,\cdot)$ names $T_a^N(p(x))$.  The same Lipschitz estimate
gives an error at most
$4^N2^{-(k+2N)}=2^{-k}$ at precision $k$.  When $x\in[0,1]$, this is
the usual nonlinear orbit $T_a^N(x)$.  No completed real-number sort or
completeness axiom occurs in the internal theorem.  Nor does the theorem
promise a fixed input precision adequate for arbitrarily many iterations:
the explicit cost is two additional bits per iteration.

\subsection{Lean status and the exact object-theory endpoint}
\label{subsec:backward-lean-status}

The file \path{lean/BackwardDemand.lean} imports only \texttt{Std} and
contains proved declarations for the following items:
\begin{itemize}
\item \texttt{horizon\_step} and \texttt{horizon\_coherence}: the schedule
  equations and their compatibility when one final step is added;
\item \texttt{compiledRecursion\_succ}: the successor equation of the
  two-loop implementation of parameter-changing recursion;
\item \texttt{transfer} and \texttt{observational\_induction}: finite
  transfer and subsequent universal closure;
\item \texttt{finite\_query\_induction}: reduction of a variable finite
  family of predecessor queries to a prefix maximum;
\item \texttt{constant\_shift\_transfer}: the explicit accumulated budget
  $K+Nc$.
\end{itemize}
All these declarations compiled with Lean 4.32.0.  Their axiom audit lists
only \texttt{propext} and \texttt{Quot.sound}; it lists neither
\texttt{sorryAx} nor \texttt{Classical.choice}.  This checks the native
Lean scheduling arguments.  The Lean predicates in these declarations
are arbitrary \texttt{Prop}-valued functions; Lean does not infer from
their types that they are primitive recursive or bounded.

The corresponding object-theory target is the repository judgment
$\texttt{Prf PRAAx}\;\Gamma\;\varphi$ \cite{repo-core}.
A complete formalization has the following concrete steps.
\begin{enumerate}
\item Supply fixed \texttt{PR} syntax for $D$, the test Booleanizing $C$,
  truncated subtraction, and the two recursions.  The recursion argument
  of \texttt{PR.prec} comes first; use projections and composition to put
  $N,K,\vec a$ in the parameter slots.  The formulas
  \eqref{eq:backward-array}--\eqref{eq:backward-horizon} are the complete
  numerical definitions needed for this step.
\item Prove the subtraction identities and
  \eqref{eq:backward-schedule-laws} as object equalities, using
  \texttt{PRA.precZero}, \texttt{PRA.precSucc}, and equality congruence.
  The schedule-coherence proof is another quantifier-free equality
  induction.  No evaluator over arbitrary \texttt{PR} syntax is required.
\item Supply the actual object derivations of bounded Booleanization and
  the application's base and local-step contracts.  For finite queries,
  use the explicit Boolean products in
  Corollary~\ref{cor:finite-query-induction} and the primitive recursive
  prefix maximum in \eqref{eq:multiple-demand-envelope}.
\item Form the syntax of \eqref{eq:backward-invariant}.  Its base has type
  $\Gamma\vdash\varphi[0/t]$.  Its step has type
  $\varphi::\operatorname{shift}(\Gamma)\vdash\varphi[S(t)/t]$.
  Use substitution, \texttt{Prf.eqSubst1}, and the supplied local contract
  to construct this step.  These shifted contexts are necessary to avoid
  capturing the arbitrary endpoint parameters.
\item Apply \texttt{PRA.qfInd} to that quantifier-free formula, then apply
  universal elimination at the term $N$.  Rewrite $h(N)=K$ and use the
  reverse Booleanization implication.  Universal introduction yields
  the corollary only when its variables are absent from the fixed context.
\item For the nonlinear application, introduce signed-rational code
  operations, clipping, and exact iteration; prove the three displayed
  rational inequalities and the parameter-shift identity.  Instantiate
  the preceding derivation transformer with their bounded tests.
\end{enumerate}
The native Lean file completes the scheduling proofs, including the
nonmonotone and variable-query cases.  The PR syntax compiler, general
Booleanization compiler, and the final object derivations of the nonlinear
example remain explicit formalization tasks.  A semantic theorem proved
with native Lean induction is not being advertised as an already
constructed \ipra\ proof object.

Finally, the rule is intentionally sufficient rather than an unrestricted
proof-extraction principle.  A local derivation and a fixed PR demand must
be supplied.  Finite consistency of a table does not prove global validity
of a proposed infinite object.  A loop that repeatedly asks whether a
represented real is nonzero has no guaranteed terminating demand merely
because every finite query is legitimate.  The diagonal obstruction to a
universal unclocked PR evaluator also remains intact
\cite{repo-contracts,repo-embedding}.
